% ---------------------------------------------------------------------------
% Chapter -- Simulation
% Work in progress: the text is being written section by section.
% Numbers come from "1 trailer/thesis_sim.m" and "2 trailers/thesis_sim.m".
% ---------------------------------------------------------------------------

\chapter{Simulation}
\label{ch:simulation}
\TODO[chapter in progress: only the outline is here]

\begin{itemize}
  \item What this chapter does: runs the controller of the previous chapters on the real kinematics, not only on the approximation.
  \item Questions to answer: does the tuned controller converge, does an adaptive gain help, what limits the period, how is the loop run in Gazebo.
\end{itemize}

% ===========================================================================
\section{Implementation}
\label{sec:sim-impl}

\begin{itemize}
  \item MATLAB files and parameters; self-test; how the bound on the controls is computed.
  \item The three simulators: approximation, exact kinematics, adaptive gain.
  \item Integration and output grid.
  \item Design used for all runs (periods, multipliers, Lyapunov functions, gains).
\end{itemize}

% ===========================================================================
\section{Closed loop on the real (non-approximated) system}
\label{sec:sim-exact}

\subsection{The design of Chapter~\ref{ch:feasibility} on the exact kinematics}

\begin{itemize}
  \item One trailer: times to 10\,\%, 1\,\% and 0.1\,\%, distance between the two models, controls, hitch angle (figures).
  \item Two trailers: first minute; comparison of the two models.
\end{itemize}

\subsection{Interval by interval over a long time}

\begin{itemize}
  \item Why the state at the end of an interval can be simulated one interval at a time.
  \item One trailer: all three curves agree.
  \item Two trailers: the norm decreases as predicted, but the limit on $u_1$ is exceeded on the exact model and the hitch angles are large.
\end{itemize}

\subsection{How far can the period be increased?}
\label{sec:sim-validity}

\begin{itemize}
  \item Deviation, progress, controls and angles against the period (table and figure).
  \item One trailer: where the contraction cap takes over.
  \item Two trailers: working limit for the period.
  \item The case of a very long period.
\end{itemize}

\subsection{From which start poses does the controller converge?}
\label{sec:sim-envelope}

\begin{itemize}
  \item Scaled starts for one trailer (table).
  \item What follows for the choice of the start and for the gain.
\end{itemize}

% ===========================================================================
\section{Fixed vs.\ adaptive gain}
\label{sec:sim-adaptive}

\begin{itemize}
  \item How the adaptive gain is computed.
  \item One trailer: time saved (figure).
  \item Two trailers: first minute.
\end{itemize}

% ===========================================================================
\section{Choice of period}
\label{sec:sim-period}

\begin{itemize}
  \item Amplitude: the period should be long.
  \item Contraction: the period should not be too long (table).
  \item Validity: the period should not be too long.
  \item Result for one and for two trailers.
\end{itemize}

% ===========================================================================
\section{Gazebo model}
\label{sec:sim-gazebo}

\begin{itemize}
  \item The model file and its two modes; passive hitches.
  \item How the state is obtained in simulation and in the laboratory.
  \item What is still open: dimensions, hitch range, response of the velocity controller.
\end{itemize}
